\usepackage{amsmath,amssymb,amsthm,mathtools}
\usepackage{graphicx,booktabs,tabularx,array}
\usepackage{algorithm,algpseudocode}
\usepackage{xcolor}
\usepackage{xurl}
\usepackage[colorlinks=true,linkcolor=blue!55!black,citecolor=blue!55!black,urlcolor=blue!55!black]{hyperref}
\newcommand{\PaperTitleEN}{Feature Geometry for Data-Pool Screening: Utility, Stability, and Limits}
\newcommand{\PaperTitleZH}{特征几何用于数据池预筛选：效用、稳定性与局限}
\newcommand{\PaperKeywordsEN}{data selection, data-pool screening, frozen representations, optimal experimental design, distribution shift, empirical evaluation}

\hypersetup{pdfauthor={Anonymous},pdfkeywords={\PaperKeywordsEN}}
\ifdefined\ChineseDraft
\hypersetup{pdftitle={\PaperTitleZH}}
\else
\hypersetup{pdftitle={\PaperTitleEN}}
\fi
\newcommand{\R}{\mathbb{R}}
\newcommand{\E}{\mathbb{E}}
\newcommand{\tr}{\operatorname{tr}}
\newcommand{\reff}{r_{\mathrm{eff}}}
\newcommand{\cC}{\mathcal{C}}
\newcommand{\cQ}{\mathcal{Q}}
\newcommand{\norm}[1]{\left\lVert#1\right\rVert}
\ifdefined\ChineseDraft
\floatname{algorithm}{算法}
\algrenewcommand{\algorithmicrequire}{输入：}
\algrenewcommand{\algorithmicensure}{输出：}
\algrenewcommand{\algorithmicfor}{遍历}
\algrenewcommand{\algorithmicdo}{执行}
\algrenewtext{EndFor}{结束遍历}
\algrenewcommand{\algorithmicreturn}{返回}
\newtheorem{proposition}{命题}
\newtheorem{lemma}{引理}
\else
\newtheorem{proposition}{Proposition}
\newtheorem{lemma}{Lemma}
\fi
\graphicspath{{figures/}}
\brokenpenalty=10000
% Generated from recorded E2b outcome tables.
\newcommand{\MainRecall}{98.333}
\newcommand{\MainRegret}{0.001311}
\newcommand{\MainReduction}{50.110}
% Supplemental diagnostics: post-hoc 2026-09-16 run.
\newcommand{\BootstrapRecall}{97.42}
\newcommand{\BootstrapMMDRecall}{97.85}
\newcommand{\BootstrapOverlap}{85.15}
\newcommand{\BrierErrorOverlap}{29.17}
\newcommand{\BrierSquaredOverlap}{30.83}
\newcommand{\BrierNLLOverlap}{90.83}
% Exploratory projection/readout follow-up, 2026-09-17.
\newcommand{\WorstErrorRecall}{40}
\newcommand{\WorstErrorGap}{1.66}
\newcommand{\ReadoutRidgeA}{98.33}
\newcommand{\ReadoutRidgeMMD}{99.38}
\newcommand{\RidgeSpan}{0.248}
\newcommand{\RidgeSelectedRegret}{0.001006}
\newcommand{\ReadoutLogisticA}{98.33}
\newcommand{\ReadoutLogisticMMD}{98.96}
\newcommand{\LogisticSpan}{14.65}
\newcommand{\LogisticSelectedRegret}{0.103}

% Record the effective layout for verify_draft.py without changing the style.
\makeatletter
\AtBeginDocument{%
  \typeout{ICLR-FORMAT-FONT: \f@size; \the\baselineskip}%
  \typeout{ICLR-FORMAT-TEXT: \the\textwidth; \the\textheight}%
  \typeout{ICLR-FORMAT-PAGE: \the\paperwidth; \the\paperheight}%
  \typeout{ICLR-FORMAT-PAR: \the\parindent; \the\parskip}%
}
\makeatother
